\subsection{豪斯多夫空间}

命题 1. 设 X 是一个拓扑空间。则下列命题等价：

\begin{enumerate}
\def\labelenumi{(\Alph{enumi})}
\setcounter{enumi}{7}
\tightlist
\item
  X 的任意两个不同的点都有不相交的邻域。
\end{enumerate}

(Hi) X 的任意一点的闭邻域之交仅由该点组成。

(Hii) 积空间 \(\mathbf{X} \times \mathbf{X}\) 的对角线是闭集。

(Hiii) 对每个集合 I，积空间 Y = X \(^{I}\) 的对角线在 Y 中是闭集。

\((\mathbf{H}^{\mathrm{iv}})\) \(\mathbf{X}\) 上没有滤子具有多于一个的极限点。

(Hv) 若 X 上的滤子 \(\mathfrak{F}\) 收敛于 \(x\)，则 \(x\) 是 \(\mathfrak{F}\) 的唯一聚点。

我们将证明蕴涵关系

以及

\[
\begin{array}{r l} (\mathrm{H}) & \Longrightarrow (\mathrm{H} ^ {\mathrm{i}}) \Longrightarrow (\mathrm{H} ^ {\mathrm{v}}) \Longrightarrow (\mathrm{H} ^ {\mathrm{iv}}) \Longrightarrow (\mathrm{H}) \\ (\mathrm{H}) & \Longrightarrow (\mathrm{H} ^ {\mathrm{iii}}) \Longrightarrow (\mathrm{H} ^ {\mathrm{ii}}) \Longrightarrow (\mathrm{H}). \end{array}
\]

\((\mathrm{H}) \Longrightarrow (\mathrm{H}^{\mathrm{i}}): \text{若 } x \neq y \text{，则存在一个开邻域 } \mathrm{U} \text{，包含 } x \text{，以及一个开邻域 } \mathrm{V} \text{，包含 } y \text{，使得 } \mathrm{U} \cap \mathrm{V} = \varnothing; \text{因此 } y \notin \overline{\mathrm{U}}.\)

\((\mathbf{H}^{\mathrm{i}}) \Longrightarrow (\mathbf{H}^{\mathrm{v}}): \text{令 } y \neq x; \text{则存在一个闭邻域 } V \text{，包含 } x \text{，使得 } y \notin V, \text{且根据假设存在 } M \in \mathfrak{F} \text{使得 } M \subset V; \text{故 } M \cap \mathbb{C}V = \emptyset. \text{但是 } \mathbb{C}V \text{是一个邻域，其中包含 } y; \text{故 } y \text{对滤子而言不是聚点，滤子为 } \mathfrak{F}.\)

\((\mathbf{H}^{\mathrm{v}}) \Rightarrow (\mathbf{H}^{\mathrm{iv}})\)：显然，因为滤子的每个极限点也是聚点。

\((\mathrm{H}^{\mathrm{iv}}) \Longrightarrow (\mathrm{H})\)：设 \(x \neq y\)，且任何邻域 \(V\)（即 \(x\) 的邻域）与任何邻域 \(W\)（即 \(y\) 的邻域）相交。于是诸集合 \(V \cap W\) 构成一个滤子的基，该滤子以 \(x\) 与 \(y\) 两者为极限点，这与假设矛盾。

\((\mathrm{H}) \Longrightarrow (\mathrm{H}^{\mathrm{iii}})\)：设 \((x) = (x_{\iota})\) 是 \(\mathbf{X}^{\mathbf{I}}\) 中不属于对角线 \(\Delta\) 的一点。则至少存在两个指标 \(\lambda, \mu\) 使 \(x_{\lambda} \neq x_{\mu}\)。设 \(\mathbf{V}_{\lambda}\)（相应地，\(\mathbf{V}_{\mu}\)）是 \(x_{\lambda}\)（相应地，\(x_{\mu}\)）在 \(\mathbf{X}\) 中的一个邻域，使得 \(\mathbf{V}_{\lambda} \cap \mathbf{V}_{\mu} = \emptyset\)；则集合 \(W = V_{\lambda} \times V_{\mu} \times \prod_{\iota \neq \lambda, \mu} X_{\iota}\)（其中 \(X_{\iota} = X\)，\(\iota \neq \lambda, \mu\)）是 \(x\) 在 \(X^{\mathbf{I}}\) 中的一个邻域（\(\S 4, no.1\)），且不与 \(\Delta\) 相交。因此 \(\Delta\) 在 \(X^{\mathbf{I}}\) 中是闭集。

\((\mathbf{H}^{\mathrm{iii}}) \Rightarrow (\mathbf{H}^{\mathrm{ii}}): \text{显然。}\)

\((H^{ii}) \Longrightarrow (H)\)：若 \(x \neq y\)，则 \((x, y) \in X \times X\) 不属于对角线 \(\Delta\)，因此（§ 4，第 1 目）存在邻域 \(V\)（即 \(x\) 的邻域）与邻域 \(W\)（即 \(y\) 的邻域），都在 \(X\) 中，使得 \((V \times W) \cap \Delta = \emptyset\)，这就意味着

\[
\mathbf {V} \cap \mathbf {W} = \varnothing .
\]

定义 1. 满足命题 1 中诸条件的拓扑空间称为豪斯多夫空间或分离空间；这种空间的拓扑称为豪斯多夫拓扑。

公理 (H) 称为豪斯多夫公理。

例. 任何离散空间都是豪斯多夫的。有理直线 Q 是豪斯多夫的，因为若 x, y 是两个有理数且 x \textless{} y，则存在有理数 z 使得 x \textless{} z \textless{} y，而 x 的邻域 {]}\(\leftarrow\) , z{[} 与 y 的邻域 {]}z, \(\rightarrow\) {[} 不相交。

一个至少含有两个点且带有最粗拓扑（§ 2，第 2 目）的集合 X 不是豪斯多夫空间。

设 \(f: X \to Y\) 是集合 \(X\) 到豪斯多夫空间 \(Y\) 中的一个映射；则由命题 1 立即得出，\(f\) 关于滤子 \(\mathfrak{F}\)（\(X\) 上的）至多有一个极限，且若 \(f\) 以 \(y\) 为关于 \(\mathfrak{F}\) 的极限，则 \(y\) 是 \(f\) 关于 \(\mathfrak{F}\) 的唯一聚点。

命题 2. 设 \(f, g\) 是拓扑空间 \(\mathbf{X}\) 到豪斯多夫空间 \(\mathbf{Y}\) 中的两个连续映射；则 \(x \in \mathbf{X}\) 且 \(f(x) = g(x)\) 的点全体组成 \(\mathbf{X}\) 中的一个闭集。

因为这个集合是 \(\mathbf{Y} \times \mathbf{Y}\) 的对角线在映射 \(x \to (f(x), g(x))\) 下的逆像，而该映射是连续的（\(\S 4\)，第 1 目，命题 1）。因此结论由 (H \(^{ii}\) ) 与 \(\S 2\)，第 1 目，定理 1 得出。

推论 1（恒等延拓原理）. 设 \(f, g\) 是拓扑空间 \(X\) 到豪斯多夫空间 \(Y\) 中的两个连续映射。若 \(f(x) = g(x)\) 在 \(X\) 的一个稠密子集的所有点处成立，则 \(f = g\)。

换句话说，X 到 Y（豪斯多夫）中的一个连续映射由它在 X 的一个稠密子集的所有点处的值唯一确定。

推论 2. 若 \(f\) 是拓扑空间 \(\mathbf{X}\) 到豪斯多夫空间 \(\mathbf{Y}\) 中的一个连续映射，则 \(f\) 的图象在 \(\mathbf{X} \times \mathbf{Y}\) 中是闭集。

因为这个图象是由所有 \((x, y) \in \mathbf{X} \times \mathbf{Y}\)（其中 \(f(x) = y\)）组成的集合，而两个映射 \((x, y) \to y\) 与 \((x, y) \to f(x)\) 都是连续的。

命题 3. 设 \((x_{i})_{1\leqslant i\leqslant n}\) 是豪斯多夫空间 \(\mathbf{X}\) 中互不相同的点组成的一个有限族；则每个 \(x_{i}\) 有一个邻域 \(\mathbf{V}_{i}\)（\(\mathbf{X}\) 中的），使得诸 \(\mathbf{V}_{i}\)（\(1\leqslant i\leqslant n\)）两两不相交。

证明对 \(n\) 用归纳法：\(n = 2\) 的情形就是公理 (H)。于是设 \(W_i\)（\(1 \leqslant i \leqslant n - 1\)）是 \(x_i\) 的邻域，使诸 \(W_i\) 两两不相交。另一方面，对 \(1 \leqslant i \leqslant n - 1\)，存在邻域 \(T_i\)（即 \(x_i\) 的邻域）与邻域 \(U_i\)（即 \(x_n\) 的邻域），二者不相交。若取 \(V_i\) 为 \(W_i \cap T_i\)（\(1 \leqslant i \leqslant n - 1\)），并取

\[
\mathrm{V} _ {n} = \bigcap_ {i = 1} ^ {n - 1} \mathrm{U} _ {i},
\]

则命题的条件即告满足。

推论. 每个有限豪斯多夫空间都是离散的。

命题 4. 豪斯多夫空间的每个有限子集都是闭集。

因为每个由单点组成的子集都依据公理 (H \(^{i}\) ) 是闭集。

命题 5. 设 X 是一个拓扑空间，并设对 X 的每对不同的点 x, y，都存在 X 到豪斯多夫空间 \(X'\) 中的一个连续映射 f，使得 \(f(x) \neq f(y)\)。则 X 是豪斯多夫的。

设 \(V'\) 与 \(W'\) 分别是 \(f(x)\) 与 \(f(y)\) 在 \(X'\) 中的不相交邻域；则 \(\overline{f}^1(V')\) 与 \(\overline{f}^1(W')\) 分别是 \(x\) 与 \(y\) 在 \(X\) 中的不相交邻域。

推论. 每个比豪斯多夫拓扑更细的拓扑都是豪斯多夫的。

